\documentclass[10pt,a4paper]{article}

%double si on veut une version prof et une version élève. simple si on veut une seule version.
\newcommand{\buildmode}{double}

\InputIfFileExists{version.tex}{}{}
% Version (définie par le script)
\providecommand{\version}{eleve}

\usepackage{feuilleactivite}

%%%à modifier à chaque fois

\renewcommand{\niveau}{1M}
\renewcommand{\chapitre}{Proportions, pourcentages}
\renewcommand{\numchap}{0.2} 
\renewcommand{\theme}{Statistiques}

\begin{document}


%\legendeicones

\vspace{-0.3cm}

\begin{objectifs}
\begin{itemize}
    \item Calculer une proportion et la mettre sous différentes formes. 
    \item Comprendre une proportion comme un coefficient de proportionalité.
    \item Retrouver un sous-effectif et un effectif total à partir d'une proportion.
    \item Calculer mentalement une proportion avec des valeurs simples.
    \item Choisir la réponse la plus vraisemblable dans un QCM.
\end{itemize}
\end{objectifs}

\prof{
    Au début du cours, bien rappeler ce qu'est une proportion, un sous-effectif et un effectif total. Rappeler les différentes formes que peuvent prendre une proportion. Faire un exemple détaillé et le laisser au tableau pendant qu'ils travaillent.
}

\section{Rappels sur les proportion et les pourcentages}
\begin{exo}[][]
Dans chacun des cas suivants, calculer la proportion et l'écrire sous forme de fraction, de nombre décimal et de pourcentage.
\begin{enumerate}
    \item Dans une classe de 25 élèves, 4 sont des gauchers.
    \item Dans un collège, 50 élèves sur 500 sont demi-pensionnaires.
    \item Dans un lycée de 100 professeurs, 9 sont professeurs de mathématiques.
\end{enumerate}
\end{exo}

\begin{exo}[][]
La saison 2022-2023 de la ligue 1 de football comptait 65\% d'entraîneurs français.
\begin{enumerate}
    \item Au total, la ligue 1 comptait 20 entraîneurs. Combien d'entraîneurs français y avait-il ?
    \item S'il y avait 60 entraîneurs au total, combien d'entraîneurs français y aurait-il ?
    \item \begin{enumerate}
        \item Complétez le tableau de proportionnalité ci-dessous.
            \begin{center}
            \begin{tabular}{|c|c|c|c|c|c|c|}
            \hline
            Nombre d'entraîneurs français &    &    &    & 65 & 130 &  \\
            \hline
            Nombre total d'entraîneurs    & 20 & 40 & 60 & 100&     & 500  \\
            \hline
            \end{tabular}
            \end{center}
        \item Combien d'entraîneurs y aurait-il si la ligue 1 comptait 130 entraîneurs français  ?
        \item Quel est le coefficient de proportionnalité du tableau ?  À quoi correspond-il dans le contexte de l'exercice ?
    \end{enumerate}
\end{enumerate} 
\end{exo}
\prof{
    Possibilité d'enlever la question 3c pour la faire à l'oral. À l'oral, faire le lien avec les élèves entre les lignes et les effectifs.\\
    À l'aide du tableau de proportionnalité, rappeler la méthode pour calculer un sous-effectif et un effecif total.
}

\begin{exo}[][]
\begin{enumerate}
    \item Dans une classe de 30 élèves, 40\% sont des filles.
    \begin{enumerate}
        \item Combien y a-t-il de filles dans la classe ?
        \item Combien y a-t-il de garçons dans la classe ?
    \end{enumerate}
    \item Dans un collège, 60\% des élèves sont demi-pensionnaires.
    \begin{enumerate}
        \item Si le collège compte 600 élèves, combien y a-t-il de demi-pensionnaires ?
        \item Si le collège compte 150 demi-pensionnaires, combien y a-t-il d'élèves au total ?
    \end{enumerate}
\end{enumerate}
\end{exo}
\prof{
    Pour la correction de la question 2, tracer un tableau de proportionnalité.
}

\begin{exo}[QCM][bac]
\begin{enumerate}
    \item Écrire le nombre \(\dfrac{3}{5}\) sous forme de pourcentage.
    \begin{tasks}(4)
        \task 0,6\%
        \task 30\%
        \task 35\%
        \task 60\%
    \end{tasks}
    \item Un lycée compte 500 élèves. 20\% d'entre eux sont externes. Le nombre d'externe est :
    \begin{tasks}(4)
        \task 10
        \task 20
        \task 100
        \task 520
    \end{tasks}
    \item Lors d'une élection, le quart des électeurs a voté pour le candidat A, 20\% a voté pour B, 34\% a voté pour C et le reste a voté pour D. Le candidat qui a eu le \underline{moins} de voix est :\begin{tasks}(4)
        \task A
        \task B
        \task C
        \task D
    \end{tasks}
\end{enumerate}
\end{exo}

\section{fréquences conditionnelles et marginales}

\begin{exo}[][]
Dans un lycée, on a répertorié l'effectif d'élève pratiquant une deuxième langue vivante en fonction de leur niveau :
\begin{center}
    \renewcommand{\arraystretch}{1.5}
    %\setlength{\tabcolsep}{20pt}
    \begin{tabular}{|l|c|c|c|c|}
        \hline
        & \textbf{LV2 espagnol} & \textbf{LV2 allemand} & \textbf{LV2 italien} & \textbf{Total} \\
        \hline
        \textbf{2nde} & 50 & 30 & 20 &  \\
        \hline
        \textbf{1ère} & 40 & 20 & 10 &  \\
        \hline
        \textbf{Terminale} & 45 & 10 & 25 & \\
        \hline
        \textbf{Total} &  & &  &  \\
        \hline
    \end{tabular}
\end{center}
\begin{enumerate}
    \item Remplir les totaux de chaque ligne et de chaque colonne.
    \medbreak
    Ce sont les \textbf{effectifs marginaux} par ligne et par colonne.
    \item Quelle est la proportion (ou fréquence) d'élèves de seconde en LV2 allemand ?
    \item Quelle est la proportion d'élèves de première en LV2 italien ?
    \medbreak
    C'est une \textbf{fréquence marginale} 
    \item Parmi les élèves en LV2 espagnol, quel est la proportion d'élèves de terminale ?
    \medbreak
    C'est une \textbf{fréquence conditionnelle} 
\end{enumerate}
\end{exo}

\begin{exo}[][]
Voici les cinq pays ayant obtenu le plus de médailles aux Jeux Olympiques de 2024 : 
\begin{center}
   % \setlength{\tabcolsep}{20pt}
\renewcommand{\arraystretch}{1.5}
\begin{tabular}{|l|c|c|c|c|}
\hline
 & \textbf{Or } & \textbf{Argent}& \textbf{Bronze} & \textbf{Total} \\
\hline
\textbf{Etats-Unis} &   &   & 42 & 126 \\
\hline
\textbf{Chine} & 40 & 27 & 24 &   \\
\hline
\textbf{Japon} & 20 & 12 & 13 & 45\\
\hline
\textbf{Australie} & 18 & 19 & 16 & 53\\
\hline
\textbf{France} & 16 &    & 22 & 64\\
\hline
\textbf{Total} & 134 & 128 & 117 & 379 \\
\hline
\end{tabular}
\end{center}
\begin{enumerate}
    \item Remplir les cases vides du tableau.  
    \item Calculer chacunes des proportions suivantes. Pour chaque proportion, préciser s'il s'agit d'une fréquence marginale ou conditionnelle.
    \begin{enumerate}
        \item La proportion de médailles obtenues par l'Australie parmi toutes les médailles d'argent.
        \item La proportion de médailles d'or obtenues par les Etats-Unis.
        \item La proportion de médailles obtenues par la France.
        \item La proportion de médailles de bronze obtenues par le Japon.
        \item La proportion de médailles d'or obtenues par la Chine parmi toutes les médailles d'or.
        \item La proportion de médailles de bronzeparmi toutes les médailles obtenues par la France.
        \item La proportion de médailles de bronze
    \end{enumerate}
\end{enumerate}
\end{exo}
\end{document}
